\subsection{链上 Möbius 算术、Smith 亏损 Euler 乘法与联合 \texorpdfstring{$\zeta$}{zeta} 的有限谱相区}\label{subsec:conclusion-chain-godel-smith-jointzeta-finite-phase}

沿用定理 \ref{thm:xi-time-part65d-chain-interior-boolean-arithmeticization}、
定理 \ref{thm:xi-chain-interior-incidence-algebra-mobius-inversion}、
定理 \ref{thm:xi-addressable-godel-worstcase-chebyshev-sharp-lower-bound}、
推论 \ref{cor:xi-chain-interior-canonical-arithmeticization-optimal-order}、
推论 \ref{cor:killo-kernel-size-general-modulus}、
定理 \ref{thm:conclusion-smith-kernel-dirichlet-euler-rational-local-factors}、
定理 \ref{thm:real-input-40-det-uv-2plus8}、
推论 \ref{cor:real-input-40-zeta-uv-sqrtv-eigs}、
定理 \ref{thm:conclusion-realinput40-uv-primitive-surgery-abel-shift}、
命题 \ref{prop:real-input-40-output-spectral-collisions} 与
注 \ref{rem:real-input-40-output-potential-degeneracy-phase} 的记号。前述链条已经分别把链上内算子格的布尔算术化、逐坐标可寻址 Gödel 外置的最坏码长主项、Smith 模核计数的 Euler 局部因子，以及真实输入 \(40\) 状态核在一维与二维输出权下的 rigid/core 因式分裂固定下来。把这些有限格、整数编码、模核亏损与联合谱分支并读之后，还可进一步压缩为下列几条终端结论：链上区间的 Möbius 函数与平方自由整除算术完全同一；按坐标可寻址的 prime-ledger 外置在整数大小 \(N\) 下至多只能承载 \((\log N)/(\log\log N)\) 量级的地址长度，而 canonical dynamic prime-register 恰以 \(T\log T\) 量级达到这一下界；Smith 信息亏损沿模数方向呈严格 Euler 乘法结构，并在最小核与零映射两端形成两极化；而在二维联合 \(\zeta\) 中，全部偶周期都被一枚与碰撞参数 \(u\) 无关的长度 \(2\) 原子统一平移，真实参数相图则只允许在有限个代数点上改换谱相。

\begin{theorem}[链上内算子格的 Möbius--算术同一化]\label{thm:conclusion-chain-interior-mobius-arithmetic-identification}
设
\[
C_n:=\{0<1<\cdots<n-1\},
\qquad
\mathcal I_n:=\operatorname{Int}(C_n),
\]
并固定两两不同素数 \(q_0,\dots,q_{n-1}\)。对任意 \(I\in\mathcal I_n\)，定义平方自由 Gödel 像
\[
\kappa(I):=\prod_{i\in \operatorname{Fix}(I)}q_i.
\]
则对任意 \(I\preceq J\) 都有
\[
\kappa(I)\mid \kappa(J),
\qquad
\operatorname{rk}[I,J]=\omega\!\left(\frac{\kappa(J)}{\kappa(I)}\right),
\]
并且区间 Möbius 函数满足
\[
\mu_{\mathcal I_n}(I,J)
=
\mu_{\mathbf N}\!\left(\frac{\kappa(J)}{\kappa(I)}\right),
\]
其中 \(\mu_{\mathbf N}\) 为经典整数 Möbius 函数，\(\omega\) 为不同素因子个数函数。特别地，\(\mathcal I_n\) 的每个闭区间都与一个平方自由整除区间规范同构。
\end{theorem}
\leanverified{paper\_conclusion\_chain\_interior\_mobius\_arithmetic\_identification}
\begin{proof}
定理 \ref{thm:xi-time-part65d-chain-interior-boolean-arithmeticization} 已证明：\(\mathcal I_n\) 是秩为 \(n-1\) 的 Boolean 格，并且 \(\kappa\) 把 \(\mathcal I_n\) 规范识别为
\[
\{d\mid Q_n:\ q_0\mid d\},
\qquad
Q_n:=\prod_{i=0}^{n-1}q_i,
\]
上的整除格。故 \(I\preceq J\) 当且仅当 \(\kappa(I)\mid\kappa(J)\)。

又由于 \(Q_n\) 平方自由，商
\[
\frac{\kappa(J)}{\kappa(I)}
\]
亦为平方自由整数，其素因子恰对应于 \(\operatorname{Fix}(J)\setminus\operatorname{Fix}(I)\)。因此
\[
\omega\!\left(\frac{\kappa(J)}{\kappa(I)}\right)
=
|\operatorname{Fix}(J)\setminus\operatorname{Fix}(I)|
=
\operatorname{rk}[I,J].
\]

最后，定理 \ref{thm:xi-chain-interior-incidence-algebra-mobius-inversion} 已给出
\[
\mu_{\mathcal I_n}(I,J)
=
(-1)^{\omega(\kappa(J)/\kappa(I))}
=
\mu_{\mathbf N}\!\left(\frac{\kappa(J)}{\kappa(I)}\right).
\]
这就得到全部结论。证毕。
\end{proof}

\begin{theorem}[可寻址素数外置编码的反向通信律]\label{thm:conclusion-addressable-prime-ledger-reverse-communication-law}
设 \(\Sigma\) 为有限字母表，\(|\Sigma|\ge 2\)。对每个 \(T\ge 1\)，考虑满足逐坐标可寻址条件的单射编码
\[
E_T:\Sigma^T\hookrightarrow \mathbf N_{\ge 1}.
\]
定义
\[
L_T:=\max_{s\in\Sigma^T}\log E_T(s),
\qquad
T_{\mathrm{addr}}(N):=
\sup\Bigl\{T:\ \exists\,E_T,\ E_T(\Sigma^T)\subseteq [1,N]\Bigr\}.
\]
再记 \(p_T\) 为第 \(T\) 个素数，\(\vartheta\) 为 Chebyshev 函数。则对任意 \(T\) 都有
\[
L_T\ge \vartheta(p_T)
=
T\log T+T\log\log T+O(T).
\]
从而
\[
T_{\mathrm{addr}}(N)\le (1+o(1))\frac{\log N}{\log\log N}
\qquad (N\to\infty).
\]
特别地，在逐坐标 prime-addressable 的素数估值编码范式中，可恢复坐标数不可能取得任何线性于 \(\log N\) 的增长律。
\end{theorem}
\begin{proof}
由定理 \ref{thm:xi-addressable-godel-worstcase-chebyshev-sharp-lower-bound}，任意满足逐坐标可寻址性的单射 \(E_T\) 都满足
\[
L_T\ge \vartheta(p_T)
=
T\log T+T\log\log T+O(T).
\]
这就给出第一式。

若某个编码 \(E_T\) 的像落在 \([1,N]\) 内，则
\[
L_T\le \log N.
\]
于是
\[
\vartheta(p_T)\le \log N.
\]
再由素数定理的等价形式
\[
\vartheta(p_T)\sim p_T\sim T\log T,
\]
对上式作渐近反演，即得
\[
T\le (1+o(1))\frac{\log N}{\log\log N}.
\]
取上确界便得到 \(T_{\mathrm{addr}}(N)\) 的上界。证毕。
\end{proof}

\begin{theorem}[canonical dynamic prime-register 的渐近最优量级]\label{thm:conclusion-canonical-dynamic-prime-register-asymptotic-optimality}
\leanverified{paper\_conclusion\_canonical\_dynamic\_prime\_register\_asymptotic\_optimality}
设 \(\Sigma\) 为有限字母表，\(|\Sigma|\ge 2\)，并取有界单射编码
\[
\operatorname{code}:\Sigma\hookrightarrow \{0,1,\dots,M\}.
\]
定义 canonical dynamic prime-register Gödel 编码
\[
E_T^{\mathrm{can}}(s_1,\dots,s_T)
:=
\prod_{t=1}^{T}p_t^{\,\operatorname{code}(s_t)},
\qquad
(s_1,\dots,s_T)\in \Sigma^T,
\]
其中 \(p_t\) 为第 \(t\) 个素数。则：
\begin{enumerate}
  \item \(E_T^{\mathrm{can}}\) 满足逐坐标可寻址性；
  \item 存在常数 \(c_1,c_2>0\) 使得对一切 \(T\ge 1\) 都有
  \[
  c_1\,T\log T
  \le
  \sup_{s\in\Sigma^T}\log E_T^{\mathrm{can}}(s)
  \le
  c_2\,T\log T;
  \]
  \item 因而，在全部逐坐标可寻址的 Gödel 外置中，\(E_T^{\mathrm{can}}\) 的最坏情形位长是渐近最优的，精确到常数因子。
\end{enumerate}
更进一步，这一最优量级可以在固定二维 \(2\)-进 ambient 中忠实实现：canonical prime-register 子单子 \(H_{\mathrm{can}}\) 通过
\[
\Psi:H_{\mathrm{can}}\hookrightarrow \operatorname{Aff}\!\bigl(T_2(E_{\min})\bigr),
\qquad
T_2(E_{\min})\cong \ZZ_2^2,
\]
得到 faithful 仿射表示。
\end{theorem}
\begin{proof}
对任意词 \(s=(s_1,\dots,s_T)\)，由算术基本定理，
\[
v_{p_t}\!\bigl(E_T^{\mathrm{can}}(s)\bigr)=\operatorname{code}(s_t)
\qquad (1\le t\le T).
\]
由于 \(\operatorname{code}\) 单射，故每个坐标 \(s_t\) 都可由第 \(t\) 个素数赋值唯一恢复，因而 \(E_T^{\mathrm{can}}\) 满足逐坐标可寻址性。

下界方面，定理 \ref{thm:conclusion-addressable-prime-ledger-reverse-communication-law} 表明，任意逐坐标可寻址编码都必须满足
\[
\max_{s\in\Sigma^T}\log E_T(s)\ge \vartheta(p_T),
\]
而 \(\vartheta(p_T)\asymp T\log T\)。故存在 \(c_1>0\) 使
\[
\sup_{s\in\Sigma^T}\log E_T^{\mathrm{can}}(s)\ge c_1\,T\log T.
\]

上界方面，由 \(\operatorname{code}(s_t)\le M\) 得
\[
\log E_T^{\mathrm{can}}(s)
=
\sum_{t=1}^{T}\operatorname{code}(s_t)\log p_t
\le
M\sum_{t=1}^{T}\log p_t
=
M\,\vartheta(p_T).
\]
再由素数定理的 Chebyshev 形式 \(\vartheta(p_T)\asymp T\log T\)，得到某个 \(c_2>0\) 使
\[
\sup_{s\in\Sigma^T}\log E_T^{\mathrm{can}}(s)\le c_2\,T\log T.
\]
两侧合并即得第 \(2\) 条，并立刻推出第 \(3\) 条。

最后，固定二维 \(2\)-进 ambient 中的 faithful realization 正是定理 \ref{thm:conclusion-prime-register-fixed-2adic-ambient-vs-finite-ledger} 的内容。证毕。
\end{proof}

\begin{theorem}[Smith 信息亏损的 Euler 乘法性与核尺度两极化]\label{thm:conclusion-smith-loss-euler-multiplicativity-polarization}
设 \(A\in M_{m\times n}(\mathbf Z)\) 满足 \(\operatorname{rank}_{\mathbf Q}(A)=m\)，其 Smith 不变量记为
\[
d_1\mid d_2\mid \cdots \mid d_m.
\]
定义模 \(b\) 的归一化亏损函数
\[
L_A(b):=\frac{\#\ker(A\bmod b)}{b^{\,n-m}}.
\]
则成立：
\begin{enumerate}
  \item \(L_A\) 是乘法函数，并且显式公式
  \[
  L_A(b)=\prod_{i=1}^m \gcd(d_i,b)
  \]
  成立；
  \item 其 Dirichlet 级数
  \[
  \mathcal D_A(s):=\sum_{b\ge1}\frac{L_A(b)}{b^s}
  \]
  在 \(\Re s>m+1\) 上绝对收敛，并具有 Euler 乘积
  \[
  \mathcal D_A(s)
  =
  \prod_{p}
  \sum_{k\ge0}
  p^{-ks+\sum_{i=1}^m \min(k,e_i(p))},
  \qquad
  e_i(p):=v_p(d_i);
  \]
  \item 存在严格的核尺度两极化：
  \[
  L_A(b)=1
  \iff
  \gcd(b,d_m)=1,
  \]
  即此时 \(\#\ker(A\bmod b)=b^{n-m}\) 取到最小值；并且
  \[
  L_A(b)=b^m
  \iff
  b\mid d_1,
  \]
  即此时 \(A\bmod b\) 为零映射，\(\#\ker(A\bmod b)=b^n\) 取到最大值。
\end{enumerate}
\end{theorem}
\begin{proof}
由推论 \ref{cor:killo-kernel-size-general-modulus}，
\[
\#\ker(A\bmod b)=b^{n-m}\prod_{i=1}^m \gcd(d_i,b),
\]
故第一条中的显式公式立即成立。

若 \(\gcd(b,c)=1\)，则
\[
\gcd(d_i,bc)=\gcd(d_i,b)\gcd(d_i,c)
\qquad (1\le i\le m),
\]
于是
\[
L_A(bc)=L_A(b)L_A(c).
\]
这证明 \(L_A\) 的乘法性。

进一步，对任意素数幂 \(b=p^k\)，有
\[
L_A(p^k)
=
\prod_{i=1}^m p^{\min(k,e_i(p))}
=
p^{\sum_{i=1}^m \min(k,e_i(p))}.
\]
因此
\[
\mathcal D_A(s)
=
\prod_p \sum_{k\ge 0} \frac{L_A(p^k)}{p^{ks}}
=
\prod_p \sum_{k\ge 0}
p^{-ks+\sum_{i=1}^m\min(k,e_i(p))}.
\]
又因 \(L_A(b)\le b^m\)，故当 \(\Re s>m+1\) 时，
\[
\sum_{b\ge1}\frac{|L_A(b)|}{|b^s|}
\le
\sum_{b\ge1}b^{m-\Re s}
<
\infty,
\]
绝对收敛成立。

对第三条，\(L_A(b)=1\) 等价于
\[
\gcd(d_i,b)=1
\qquad (1\le i\le m).
\]
由于 \(d_i\mid d_m\)，这又等价于 \(\gcd(d_m,b)=1\)。故最小核恰在 \(b\) 与最大 Smith 不变量互素时出现。

另一方面，\(L_A(b)=b^m\) 等价于
\[
\gcd(d_i,b)=b
\qquad (1\le i\le m),
\]
即 \(b\mid d_i\) 对所有 \(i\) 同时成立。取 \(i=1\) 即得 \(b\mid d_1\)；反过来，由 \(d_1\mid d_i\) 可知 \(b\mid d_1\) 蕴含 \(b\mid d_i\) 对全部 \(i\) 成立。此时 Smith 形
\[
\operatorname{diag}(d_1,\dots,d_m,0,\dots,0)
\]
模 \(b\) 后退化为零矩阵，因此 \(A\bmod b\) 为零映射，核大小为 \(b^n\)。证毕。
\end{proof}

\begin{theorem}[二维联合 \texorpdfstring{$\zeta$}{zeta} 的普适长度二原子与奇偶迹分裂]\label{thm:conclusion-realinput40-joint-zeta-lengthtwo-atom-trace-splitting}
在真实输入 \(40\) 状态核的二维联合加权下，记
\[
\Delta(z;u,v)=\det(I-zM_{u,v}),
\qquad
\zeta(z;u,v)=\Delta(z;u,v)^{-1}.
\]
由定理 \ref{thm:real-input-40-det-uv-2plus8}，
\[
\Delta(z;u,v)=(1-vz^2)\Delta_{\mathrm{core}}(z;u,v),
\qquad
\Delta_{\mathrm{core}}(z;u,v):=-P_8(z;u,v),
\]
且 \(\Delta_{\mathrm{core}}(0;u,v)=1\)。定义
\[
\zeta_{\mathrm{core}}(z;u,v):=\Delta_{\mathrm{core}}(z;u,v)^{-1}.
\]
则有精确分解
\[
\zeta(z;u,v)=\frac{1}{1-vz^2}\,\zeta_{\mathrm{core}}(z;u,v).
\]
若再以
\[
\log \zeta_{\mathrm{core}}(z;u,v)
=
\sum_{n\ge1}\frac{\tau_n^{\mathrm{core}}(u,v)}{n}z^n
\]
定义 core 迹系数，则对每个 \(r\ge 1\) 都有
\[
\operatorname{Tr}(M_{u,v}^{2r-1})=\tau_{2r-1}^{\mathrm{core}}(u,v),
\]
\[
\operatorname{Tr}(M_{u,v}^{2r})=2v^r+\tau_{2r}^{\mathrm{core}}(u,v).
\]
换言之，奇周期统计完全由八次 core 因子承担，而全部偶周期都被一枚与碰撞参数 \(u\) 无关的刚性长度 \(2\) 原子统一平移。
\end{theorem}
\begin{proof}
由定理 \ref{thm:real-input-40-det-uv-2plus8}，\(\Delta(z;u,v)\) 具有精确因式分解
\[
\Delta(z;u,v)=(1-vz^2)\Delta_{\mathrm{core}}(z;u,v),
\]
故
\[
\zeta(z;u,v)=\frac{1}{1-vz^2}\,\zeta_{\mathrm{core}}(z;u,v).
\]

对两端取对数，得到
\[
\log\zeta(z;u,v)
=
-\log(1-vz^2)+\log\zeta_{\mathrm{core}}(z;u,v)
=
\sum_{r\ge1}\frac{v^r}{r}z^{2r}
\;+\;
\sum_{n\ge1}\frac{\tau_n^{\mathrm{core}}(u,v)}{n}z^n.
\]
另一方面，由 Fredholm 行列式的标准迹展开，
\[
\log\zeta(z;u,v)
=
\sum_{n\ge1}\frac{\operatorname{Tr}(M_{u,v}^n)}{n}z^n.
\]
逐项比较系数即可得到：奇数次项不受 \(-\log(1-vz^2)\) 影响，从而
\[
\operatorname{Tr}(M_{u,v}^{2r-1})=\tau_{2r-1}^{\mathrm{core}}(u,v),
\]
而偶数次项满足
\[
\frac{\operatorname{Tr}(M_{u,v}^{2r})}{2r}
=
\frac{v^r}{r}+\frac{\tau_{2r}^{\mathrm{core}}(u,v)}{2r},
\]
等价于
\[
\operatorname{Tr}(M_{u,v}^{2r})=2v^r+\tau_{2r}^{\mathrm{core}}(u,v).
\]
证毕。
\end{proof}

\begin{theorem}[正实参数轴的有限谱相剖分]\label{thm:conclusion-realinput40-positive-axis-finite-spectral-phase-partition}
设 \(Q(z,u)\) 为命题 \ref{prop:real-input-40-output-spectral-collisions} 中从
\[
\Delta(z,u)=(1-uz^2)Q(z,u)
\]
分离出的 core 因子，并令
\[
u_1<u_2<u_3<u_4<u_5<u_6
\]
为次数 \(14\) 的不可约多项式 \(D(u)\) 在正实轴上的六个根，即
\[
\begin{aligned}
u_1&\approx 0.0084602654608788,\\
u_2&\approx 0.7455344958029896,\\
u_3&\approx 0.9328372624859649,\\
u_4&\approx 1.0547964213679192,\\
u_5&\approx 2.0903923780490870,\\
u_6&\approx 15.5228207044268520.
\end{aligned}
\]
则正实参数轴
\[
(0,\infty)\setminus\{u_1,u_2,u_3,1,u_4,u_5,u_6\}
\]
恰分解为八个开区间；并且在每个开区间上，\(Q(z,u)\) 的全部实谱支都保持单根，其实根排序与次模态类型保持不变。因而该一参数实谱族在 \(u>0\) 上至多存在八个实谱相区，而每一次相变都被上述七个代数点严格钉扎。
\end{theorem}
\begin{proof}
命题 \ref{prop:real-input-40-output-spectral-collisions} 已给出
\[
\operatorname{Res}_z(Q,Q_z)=4u^{15}(u-1)^3D(u),
\]
并指出：在 \(u>0\) 上，除去结构性退化点 \(u=1\) 之外，\(Q(z,u)\) 出现重根当且仅当 \(D(u)=0\)。

另一方面，注 \ref{rem:real-input-40-output-potential-degeneracy-phase} 已列出 \(D(u)\) 的六个正根
\[
u_1<u_2<u_3<u_4<u_5<u_6,
\]
并说明正实 \(u\)-轴正是被这些点与 \(u=1\) 划分为若干开区间；在每个区间内，实谱支的排序与次模态类型保持恒定，不会发生谱支碰撞或交换。

因此，对集合
\[
\Sigma:=\{u_1,u_2,u_3,1,u_4,u_5,u_6\}
\]
而言，\((0,\infty)\setminus\Sigma\) 的各个连通分支恰是全部可能的实谱相区。由于 \(|\Sigma|=7\)，故其开区间总数正好为 \(8\)。证毕。
\end{proof}

\begin{theorem}[稳定 \texorpdfstring{$K$}{K}-骨架与瞬态 Jordan 塔的严格分离]\label{thm:conclusion-realinput40-bf-jordan-tower-separation}
设 \(A=M_{\mathrm{ess}}\) 为真实输入 \(40\) 状态核的 essential \(20\times 20\) 核矩阵，\(M\) 为 full \(40\times 40\) 核矩阵。则二者具有完全相同的 Bowen--Franks 稳定骨架
\[
\mathrm{BF}(A)\cong \mathrm{BF}(M)\cong \ZZ^2\oplus \ZZ/2\ZZ,
\]
但零谱幂零深度严格不同：
\[
r_{\mathrm{ess}}=3,\qquad r_{\mathrm{full}}=5.
\]
等价地，full 核相对于 essential 核并未改变任何 Bowen--Franks 型稳定不变量，却精确增加了两层纯瞬态的 \(0\)-谱 Jordan 阴影。
\end{theorem}
\begin{proof}
命题 \ref{prop:real-input-40-bf-ktheory} 已给出
\[
\mathrm{BF}(A)\cong \ZZ^2\oplus \ZZ/2\ZZ.
\]
而注 \ref{rem:real-input-40-bf-ktheory-repro} 连同其 Smith 标准形审计表明，full 矩阵 \(M\) 的 \(I-M\) 也导出同一 Bowen--Franks 群，故
\[
\mathrm{BF}(A)\cong \mathrm{BF}(M)\cong \ZZ^2\oplus \ZZ/2\ZZ.
\]

另一方面，命题 \ref{prop:real-input-40-nilpotent-index} 已证明
\[
\dim\ker A^k=(7,10,11)\ \text{并于 }k=3\text{ 稳定},
\qquad
\dim\ker M^k=(16,24,29,30,31)\ \text{并于 }k=5\text{ 稳定}.
\]
因此 \(0\) 处最大 Jordan 块大小分别为
\[
r_{\mathrm{ess}}=3,\qquad r_{\mathrm{full}}=5.
\]
稳定骨架相同而幂零深度增加 \(2\)，即得结论。证毕。
\end{proof}

\begin{corollary}[full \texorpdfstring{$\to$}{to} essential 压缩只是纯零谱瞬态剥离]\label{cor:conclusion-realinput40-full-essential-pure-zero-spectrum-stripping}
真实输入核从 full \(40\) 状态表示压缩到 essential \(20\) 状态核心时，所失去的并不是任何新的稳定 \(K\)-理论信息，而只是两层附着在 \(0\)-谱上的瞬态 Jordan 塔。
\end{corollary}
\begin{proof}
定理 \ref{thm:conclusion-realinput40-bf-jordan-tower-separation} 已表明二者的 Bowen--Franks 骨架完全一致，而零谱幂零深度由 \(5\) 降到 \(3\)。因此压缩所剥离者只能是纯 \(0\)-谱瞬态层。证毕。
\end{proof}

\begin{theorem}[同步长度与零谱湮灭时标的完全对齐]\label{thm:conclusion-realinput40-sync-length-zero-spectrum-alignment}
\leanverified{paper\_conclusion\_realinput40\_sync\_length\_zero\_spectrum\_alignment}
对 full 核 \(M\)，零谱瞬态的全部可观测影响在且仅在步长
\[
m\ge 5
\]
后彻底湮灭；而这一临界时标恰与最短同步复位词长度
\[
|0^5|=5
\]
完全一致。对 essential 核 \(A\)，相应阈值则为
\[
m\ge 3.
\]
因而 full 核相对于 essential 核多出的两层 Jordan 阴影，正好对应于同步建立前的两步预核心过渡层。
\end{theorem}
\begin{proof}
由推论 \ref{cor:real-input-40-nilpotent-memory-depth}，当 \(m\ge r_{\mathrm{full}}=5\) 时，一切形如
\[
f_m=v^\ast M^m w
\]
的可观测量都与 \(w\) 在 \(V_0(M)\) 上的分量无关；同理，essential 核的对应阈值为 \(r_{\mathrm{ess}}=3\)。

反过来，由命题 \ref{prop:real-input-40-nilpotent-index} 的 Jordan 指数尖锐性，存在 \(w_0\in V_0(M)\) 使 \(M^4w_0\neq 0\)。取某个线性泛函 \(v^\ast\) 使 \(v^\ast M^4w_0\neq 0\)，便知在 \(m=4\) 时零谱瞬态贡献仍可被观测；essential 情形同理，在 \(m=2\) 时仍可留下非零贡献。故两个阈值都是尖锐的。

另一方面，命题 \ref{prop:real-input-40-reset-word} 已给出：\(0^5\) 是真实输入核的最短同步复位词，长度恰为 \(5\)。因此 full 核零谱瞬态的完全湮灭时标与最短同步长度完全一致，而 essential 阈值较之少 \(2\) 步，正对应于 full 相对 essential 额外携带的两层瞬态塔。证毕。
\end{proof}

\begin{corollary}[同步是零谱瞬态被完全抹去的精确时间边界]\label{cor:conclusion-realinput40-sync-exact-zero-spectrum-boundary}
在真实输入 \(40\) 状态核中，“同步建立”并非附加语义，而是 \(0\)-谱瞬态对所有有限维观测完全失效的精确时间边界。
\end{corollary}
\begin{proof}
由定理 \ref{thm:conclusion-realinput40-sync-length-zero-spectrum-alignment}，full 核的同步临界长度与零谱湮灭阈值同为 \(5\)。两者完全重合，故同步就是零谱瞬态被完全抹去的精确边界。证毕。
\end{proof}

\begin{theorem}[Parry 基线下的复位门湮灭--再生定理]\label{thm:conclusion-realinput40-reset-gate-annihilation-regeneration}
记
\[
Z_{\ge 5}:=\{d_{k-4}\cdots d_k=0^5\}
\]
为五零复位门事件。则在 Parry 最大熵基线下：
\begin{enumerate}
  \item 每次 \(Z_{\ge 5}\) 发生时，后继同步态被强制压到唯一状态 \(000\)，从而完全遗忘先前同步不确定性；
  \item 同时，所有来自 \(V_0(M)\) 的零谱瞬态贡献在该门后被彻底湮灭；
  \item 此类完全遗忘门以显式闭式强度反复出现：
  \[
  \mathbb P(Z_{\ge 5})=\frac{9-4\sqrt5}{5},
  \qquad
  \mathbb E[\tau_{Z_{\ge 5}}]=45+20\sqrt5,
  \]
  并且从典型非复位时刻到下一次复位门的条件等待时间满足
  \[
  \mathbb E[T]=\frac{555}{8}+\frac{127}{4}\sqrt5.
  \]
\end{enumerate}
\end{theorem}
\begin{proof}
第一条正是命题 \ref{prop:real-input-40-reset-word} 的同步复位性质：
\[
\delta(s,0^5)=000\qquad(\forall s\in Q).
\]
第二条由定理 \ref{thm:conclusion-realinput40-sync-length-zero-spectrum-alignment} 立即得出，因为复位门长度恰等于 full 核零谱瞬态的完全湮灭阈值 \(5\)。

第三条则由命题 \ref{prop:real-input-40-reset-regeneration-constants} 给出：
\[
\mathbb P(Z_{\ge 5})=\frac{9-4\sqrt5}{5},
\qquad
\mathbb E[\tau_{Z_{\ge 5}}]=45+20\sqrt5,
\qquad
\mathbb E[T]=\frac{555}{8}+\frac{127}{4}\sqrt5.
\]
故 \(Z_{\ge 5}\) 不仅是同步复位词，而且是携带完全闭式再生强度的瞬态湮灭门。证毕。
\end{proof}

\begin{corollary}[同步与零谱的双重湮灭门]\label{cor:conclusion-realinput40-sync-zero-spectrum-double-gate}
真实输入核中的 \(0^5\) 不是近似再生事件，而是一个严格的同步--零谱双重湮灭门：它一方面把内部状态压到唯一同步态，另一方面把全部 \(0\)-谱瞬态贡献同时归零。
\end{corollary}
\begin{proof}
由定理 \ref{thm:conclusion-realinput40-reset-gate-annihilation-regeneration} 的前两条，二者在同一个五零复位门上同时发生。证毕。
\end{proof}

\begin{theorem}[合法性记忆层对周期谱量的完全失语]\label{thm:conclusion-realinput40-legality-memory-spectral-aphasia}
\leanverified{paper\_conclusion\_realinput40\_legality\_memory\_spectral\_aphasia}
在一阶确定性在线归一化类别中，真实输入约束的精确实现至少需要 \(40\) 个状态；但与一切周期谱量相关的非零动力学却已经完全压缩到唯一的 essential \(20\) 状态核心 \(A=M_{\mathrm{ess}}\) 上。更具体地，
\[
\det(I-zM)=\det(I-zA)
=(1-z)^2(1+z)^3(1-3z+z^2)(1+z-z^2),
\]
故 full \(40\) 状态核与 essential \(20\) 状态核具有完全相同的非零谱
\[
\{1,1,-1,-1,-1,\varphi^2,\varphi^{-2},-\varphi,\varphi^{-1}\},
\]
并因而诱导相同的 Dirichlet 化 \(\zeta\)、primitive 轨道计数、Abel--Mertens 常数以及一切只锚定非零谱的周期不变量。换言之，额外的 \(20\) 个合法性记忆状态在实现上不可删，但对全部周期谱量完全失语。
\end{theorem}
\begin{proof}
定理 \ref{thm:real-input-40-product-lower-bound} 已证明：若要求一阶确定性在线归一化与输入合法性同时显式实现，则状态数下界恰为
\[
|\mathrm{States}|\ge 40.
\]
另一方面，命题 \ref{prop:real-input-40-essential-20} 给出 full 核中的唯一 essential 强连通分量恰有 \(20\) 个状态；而定理 \ref{thm:conclusion-realinput40-bf-jordan-tower-separation} 与推论 \ref{cor:conclusion-realinput40-full-essential-pure-zero-spectrum-stripping} 已经说明，full 核相对于 essential 核新增的部分只是一座纯 \(0\)-谱瞬态 Jordan 塔，不改变任何稳定的非零谱骨架。

再由命题 \ref{prop:real-input-40-fib-tensor} 及其推论 \ref{cor:real-input-40-spectrum-decomp}，
\[
\det(I-zM)=\det(I-zA)
=(1-z)^2(1+z)^3(1-3z+z^2)(1+z-z^2),
\]
故 full 核与 essential 核的非零谱完全一致，且恰为
\[
\{1,1,-1,-1,-1,\varphi^2,\varphi^{-2},-\varphi,\varphi^{-1}\}.
\]
于是凡仅由 \(\det(I-zM)\)、\(\zeta_M\)、\(\Tr(M^n)\)、primitive Möbius 反演与相应 Abel 有限部决定的对象，全部都只能看见 essential 核而看不见附加的合法性记忆层。故这额外的 \(20\) 个状态虽然对忠实实现不可省略，却对全部周期谱量完全静默。证毕。
\end{proof}

\endinput


\noindent
上述链条把若干原本分散的有限机制压缩为同一幅图景：离散侧，链上内算子格与平方自由整除格完全同一，而逐坐标可寻址的 prime-ledger 外置在整数大小 \(N\) 下至多只能承载 \((\log N)/(\log\log N)\) 量级的地址长度；与此同时，canonical dynamic prime-register 又以 \(T\log T\) 的最坏码长达到这一正确量级，并能在固定二维 \(2\)-进 Tate ambient 中忠实实现。模算术侧，Smith 信息亏损沿模数方向呈严格 Euler 乘法结构，并在最小核与零映射两端形成刚性的两极化；解析谱侧，二维联合 \(\zeta\) 由一枚普适长度 \(2\) 原子与八次 core 因子精确分裂，偶周期统一获得刚性平移，而真实参数相图又只允许在有限个代数点上改换谱相。再把真实输入 \(40\) 状态核的 Smith/Bowen--Franks 骨架、\(0\)-谱 Jordan 塔与同步复位门并读，则同一 finite-phase 机制还呈现出更尖锐的三层分裂：稳定 \(K\)-骨架保持不变，full 相对 essential 只多出两层同步前的瞬态阴影，而 \(0^5\) 则同时充当同步态与零谱贡献的双重湮灭门。因而，本文中的“算术化”“亏损”“谱分支”与“同步复位”并不平行，而是共同落在一个由平方自由整除、Euler 乘积、有限代数碰撞、动态 prime-ledger 最优量级与幂零瞬态剥离所支配的统一有限机制之中。

\endinput
